This research investigates whether the \newterm{\fv} framework, which has been successful in representing simple in-context learning tasks, can generalize to complex generative tasks like abstractive summarization. We hypothesize that these vectors can decode the ``default constraints''---such as \extractive versus \abstractive styles---that Large Language Models (\fvsabbr) employ by default. We extract layer-wise task representations from \gpttwo during few-shot summarization across two datasets with opposing stylistic biases: \cnndm (\extractive) and \xsum (\abstractive). Our analysis reveals a highly unified internal representation for the summarization task across early and middle layers (cosine similarity $> 0.93$), which sharply diverges only at the final output layer (similarity drops to $0.39$). These findings suggest that \fvsabbr treat summarization as a single, generalizable task internally, applying specific stylistic constraints as a late-stage linear shift before vocabulary projection. Our work advances representation engineering by demonstrating that complex, multi-constraint tasks can be decoded and potentially controlled through targeted activation interventions in the final layers of the network.
